% Conclusao (todos). Rascunho escrito a partir dos resultados; revisar em grupo.
\section{Conclusão}

Este trabalho implementou duas estratégias para identificar pontes em grafos simples não-direcionados, o método ingênuo e o método de Tarjan~\cite{tarjan1974}, e as usou dentro do algoritmo de Fleury para encontrar caminhos eulerianos. As duas estratégias implementam a mesma interface, de modo que o Fleury é o mesmo código nos dois casos e a comparação mede apenas a diferença entre elas. Os algoritmos foram avaliados em grafos aleatórios conexos e esparsos, eulerianos, semi-eulerianos e não eulerianos, com até 100.000 vértices, e todas as execuções terminaram dentro do limite de 30 minutos com resultado correto.

Os experimentos confirmaram a análise de complexidade. Para encontrar todas as pontes de um grafo fixo, o método de Tarjan cresce de forma linear e foi quase 6.000 vezes mais rápido que o método ingênuo com 100.000 vértices. Dentro do Fleury, porém, a vantagem se inverte: como o grafo muda a cada passo, cada teste de ponte exige uma nova busca nas duas estratégias, e ambas custam $O(E \cdot (V + E))$. Nesse cenário, o teste ingênuo, que interrompe a busca assim que encontra a outra extremidade da aresta, foi de 2,8 a 7,5 vezes mais rápido que o método de Tarjan, que percorre sempre a componente inteira. Nos grafos não eulerianos, o Fleury termina na verificação inicial de paridade e conectividade, em tempo linear e igual para as duas estratégias.

Assim, a escolha da estratégia depende do uso: o método de Tarjan é a melhor opção para identificar todas as pontes de uma vez, mas não melhora o Fleury quando aplicado a cada passo. Como trabalho futuro, o Fleury poderia usar estruturas que mantêm as pontes atualizadas conforme as arestas são removidas, evitando refazer a busca a cada passo. Outra alternativa é o algoritmo de Hierholzer, que encontra o caminho euleriano em tempo $O(V + E)$ sem precisar identificar pontes.
